\documentclass[12pt, twoside]{article}
\input{../preamble}

\fancyhead[LE]{\thepage}
\fancyhead[RO]{Name: \hspace{4cm} \,\\ \hspace{2.9cm} \,}
\fancyhead[LO]{La Scuola d'Italia / Dr.\ Huson / IB Math \\* 7 March 2026}

\begin{document}
\subsubsection*{5.1 Classwork: Introduction to Kinematics}

\textbf{Key definitions.}
A \emph{particle} moves along a straight line.
Its \textbf{position} $s$ is measured from a fixed origin, with positive to the right.

\smallskip
\begin{tabular}{r@{\;---\;}l@{\qquad}l}
  \textbf{Displacement} & change in position & \emph{vector} (has sign: direction matters) \\
  \textbf{Distance}     & total length of path & \emph{scalar} (always $\geq 0$) \\
  \textbf{Velocity}     & rate of change of displacement &
    $v = \dfrac{\Delta s}{\Delta t}$ \;(vector) \\[4pt]
  \textbf{Speed}        & magnitude of velocity & $|v|$ \;(scalar, always $\geq 0$) \\
  \textbf{Acceleration} & rate of change of velocity &
    $a = \dfrac{\Delta v}{\Delta t}$ \;(vector)
\end{tabular}

\medskip

\begin{enumerate}[itemsep=0.25cm]

%--- Problem 1: displacement vs distance ---
\item A runner jogs 400\,m east and then 150\,m west.
\begin{enumerate}[itemsep=0.15cm]
  \item Find the total \textbf{distance} the runner travels. \hfill [1]
  \item Find the runner's \textbf{displacement} from the starting point.
        State the direction. \hfill [2]
  \item Explain why these two answers are different. \hfill [1]
\end{enumerate}

%--- Problem 2: velocity from position-time graph ---
\item The graph below shows the position $s$ (metres) of a particle
      at time $t$ (seconds).

\begin{center}
\vspace{-0.3cm}
\begin{tikzpicture}[scale=0.6]
  \draw[-{Stealth}] (-0.5,0) -- (9,0) node[right] {$t$ (s)};
  \draw[-{Stealth}] (0,-2.5) -- (0,5.5) node[above] {$s$ (m)};
  \foreach \x in {1,...,8} \draw[gray!30] (\x,-2.5)--(\x,5.5);
  \foreach \y in {-2,...,5} \draw[gray!30] (-0.5,\y)--(9,\y);
  \foreach \x in {1,2,...,8} \draw (\x,0.1)--(\x,-0.1) node[below]{\small\x};
  \foreach \y in {-2,-1,1,2,3,4,5} \draw (0.1,\y)--(-0.1,\y) node[left]{\small\y};
  \node[below left] at (0,0) {\small 0};
  \draw[very thick, blue]
    (0,0) -- (2,4) -- (5,4) -- (7,-2) -- (8,-2);
  \foreach \pt in {(0,0),(2,4),(5,4),(7,-2),(8,-2)}
    \fill[blue] \pt circle (2.5pt);
\end{tikzpicture}
\vspace{-0.3cm}
\end{center}

\begin{enumerate}[itemsep=0.15cm]
  \item Find the velocity of the particle during $0 \leq t \leq 2$. \hfill [2]
  \item Describe the motion of the particle during $2 \leq t \leq 5$. \hfill [1]
  \item Find the velocity during $5 \leq t \leq 7$. \hfill [2]
  \item Find the total \textbf{distance} travelled from $t=0$ to $t=7$. \hfill [2]
  \item Find the \textbf{displacement} from $t=0$ to $t=7$. \hfill [1]
\end{enumerate}

\newpage

%--- Problem 3: velocity-time graph, acceleration ---
\item The velocity $v$ (m\,s$^{-1}$) of a particle is shown below.

\begin{center}
\vspace{-0.3cm}
\begin{tikzpicture}[scale=0.6]
  \draw[-{Stealth}] (-0.5,0) -- (9,0) node[right] {$t$ (s)};
  \draw[-{Stealth}] (0,-1.5) -- (0,5.5) node[above] {$v$ (m\,s$^{-1}$)};
  \foreach \x in {1,...,8} \draw[gray!30] (\x,-1.5)--(\x,5.5);
  \foreach \y in {-1,...,5} \draw[gray!30] (-0.5,\y)--(9,\y);
  \foreach \x in {1,2,...,8} \draw (\x,0.1)--(\x,-0.1) node[below]{\small\x};
  \foreach \y in {-1,1,2,3,4,5} \draw (0.1,\y)--(-0.1,\y) node[left]{\small\y};
  \node[below left] at (0,0) {\small 0};
  \draw[very thick, red]
    (0,0) -- (3,3) -- (6,3) -- (8,0);
  \foreach \pt in {(0,0),(3,3),(6,3),(8,0)}
    \fill[red] \pt circle (2.5pt);
\end{tikzpicture}
\vspace{-0.3cm}
\end{center}

\begin{enumerate}[itemsep=0.15cm]
  \item Find the acceleration during $0 \leq t \leq 3$. \hfill [2]
  \item State the acceleration during $3 \leq t \leq 6$. \hfill [1]
  \item Find the acceleration during $6 \leq t \leq 8$. \hfill [2]
  \item Is the particle slowing down during $6 \leq t \leq 8$?
        Explain your answer. \hfill [1]
\end{enumerate}

\medskip

%--- Problem 4: scalar vs vector identification ---
\item For each quantity, state whether it is a \textbf{scalar} or a \textbf{vector}. \hfill [4]
\begin{multicols}{2}
\begin{enumerate}[itemsep=0.1cm]
  \item A car travels at 60\,km/h.
  \item A car drives 3\,km north.
  \item Velocity $v = -5$\,m\,s$^{-1}$.
  \item Temperature is $22^\circ$C.
\end{enumerate}
\end{multicols}

%--- Problem 5: negative velocity interpretation ---
\item A particle moves along a line. At time $t=4$\,s, its velocity is $v = -3$\,m\,s$^{-1}$.
\begin{enumerate}[itemsep=0.15cm]
  \item What does the negative sign tell you about the particle's motion? \hfill [1]
  \item Find the speed of the particle at $t=4$. \hfill [1]
\end{enumerate}

%--- Problem 6: calculate velocity and acceleration from table ---
\item The table shows the position of a particle at different times.

\renewcommand{\arraystretch}{1.3}
\begin{center}
\begin{tabular}{|c|c|c|c|c|c|}
  \hline
  $t$ (s) & 0 & 1 & 2 & 3 & 4 \\
  \hline
  $s$ (m) & 2 & 5 & 8 & 11 & 14 \\
  \hline
\end{tabular}
\end{center}
\vspace{-0.2cm}

\begin{enumerate}[itemsep=0.15cm]
  \item Find the velocity of the particle. \hfill [2]
  \item Find the acceleration of the particle. \hfill [1]
  \item Describe the motion in words. \hfill [1]
\end{enumerate}

%--- Problem 7: synthesis ---
\item A cyclist starts from rest and accelerates uniformly at
      $2$\,m\,s$^{-2}$ for 5 seconds.
\begin{enumerate}[itemsep=0.15cm]
  \item Find the velocity of the cyclist at $t=5$\,s. \hfill [2]
  \item Sketch a velocity--time graph for $0 \leq t \leq 5$. \hfill [2]
\end{enumerate}

\end{enumerate}
\end{document}
